\section{Getting started in AI safety}
\label{sec:conclusion}

Technological progress is a human choice, not a law of nature. Mathematics holds truth-seeking as a core value, but the civic mathematician also asks what our truths are for, and whether they can help design AI that is more legible, steerable, and cooperative with us.

The open problems in this invitation are meant as entry points into the new field of Math for AI Safety. Many more problems can be found in the \href{https://github.com/lionellevine/MAIS/blob/main/open-problems/README.md}{MAIS repository}, and the \href{https://mbrcic.github.io/ai-safety-formalization-atlas/}{AI Safety Formalization Atlas}, maintained by Mario Br\v{c}i\'c, collects machine-checked Lean proofs of results relevant to AI safety, including verdicts on submitted solutions to MAIS problems. I hope you'll pick a problem whose language already feels like home, strip it down to the simplest interesting case, and start proving things!

The rest of this section collects useful links to the growing ecosystem of AI safety organizations and funders.

\subsection*{Upskilling}
The \href{https://www.arena.education/}{ARENA} curriculum teaches the engineering side; the month-long \href{https://www.iliad.ac/intensive}{Iliad Intensive} is a full-time course on the theory, aimed at mathematicians, physicists, and theoretical computer scientists. \href{https://www.matsprogram.org/}{MATS} and \href{https://sparai.org/}{SPAR} pair newcomers with mentors for a first research project.

\subsection*{Institutes}
Research organizations in AI safety include \href{https://www.alignment.org/}{ARC}, \href{https://humancompatible.ai/}{CHAI}, \href{https://constellation.org/}{Constellation}, \href{https://www.far.ai/}{FAR AI}, \href{https://maisi.org/}{MAISI}, \href{https://resi.org/}{RESI}, and \href{https://resolution.org}{Resolution}. Many run long-term fellowship or visitor programs.

\subsection*{Meetings}
\href{https://www.iliadconference.com}{ILIAD} is a conference on mathematical approaches to AI alignment, and FAR AI runs a series of \href{https://www.far.ai/events}{Alignment Workshops}.

\subsection*{Funding}
Grantmakers supporting work in AI safety include \href{https://coefficientgiving.org/funds/navigating-transformative-ai/}{Coefficient Giving}, \href{https://www.cooperativeai.com/}{CAIF}, \href{https://survivalandflourishing.fund/}{SFF}, and \href{https://alignmentproject.aisi.gov.uk/}{UK AISI}.

\section*{Acknowledgments}

I thank Jesse Hoogland, Hyojeong Son, Jacob Tsimerman, Claude, and Codex for many inspiring conversations. Scott Aaronson, Ahmed Bou-Rabee, Vince Conitzer, Jonathan Gabor, Chris Hillar, Brad Knox, Phil Sosoe, Samuel Speas, Kate Stange, Steve Strogatz, Ariel Yadin, and an anonymous referee provided valuable feedback on an early draft.

\section*{AI collaborators}
\label{sec:ai-collaborators}

I wrote the first draft in collaboration with Claude Opus 4. To start, I instructed Opus to read the full text of approximately 100 AI safety papers and produce a summary of the main results and techniques of each, along with possible relevance to this invitation paper. I supplied the overarching structure of the invitation (organized by field, ending each section with an open problem) and wrote parts of the introduction. Opus then used its summaries along with samples of my writing to write a draft of each section in my voice, which I edited extensively. Later revisions were polished with the help of Claude Fable 5, GPT 5.6 Sol, and GPT 6 Astra.

In July 2026 an exuberant Fable produced approximately 200 pages worth of open problems in a single night while I was asleep. I asked Sol to audit these for openness, well-posedness, and plausible AI safety relevance. The problems that passed this audit were used to seed the \href{https://github.com/lionellevine/MAIS/blob/main/open-problems/README.md}{master list of open problems} in the \href{https://github.com/lionellevine/MAIS}{Math for AI Safety (MAIS) repository}, a new hub for open-source collaboration. 

The audit also turned up a small surprise: while checking the problem that became \texttt{MAIS-O1}, Sol found a gap in the published proof of Critch's bounded L\"ob theorem, and proposed a repair (\href{https://github.com/lionellevine/MAIS/tree/main/papers/P2}{MAIS-P2}). Critch corrected the hypothesis (\href{https://x.com/AndrewCritchPhD/status/2082141000831795467}{post on X, July 28, 2026}) and a version of his theorem, in a bespoke proof calculus, has been proved in Lean~\cite{duclaux2026}.

\appendix

\section{Glossary of machine learning terms for mathematicians}
\label{app:glossary}

\begin{description}
\item[\glsentry{neural-network}{neural network}] a function with many (up to trillions of) adjustable real parameters (its \emph{weights}), built by composing simple layers and tuned by \emph{gradient descent} to reduce a training \emph{loss}.
\item[\glsentry{layer}{layer}] a map $x\mapsto\sigma(Wx+b)$: an affine map, whose matrix and vector entries are weights, followed by a fixed nonlinearity $\sigma$ applied to each coordinate; a network is a composition of layers (a transformer alternates them with attention layers, which mix across positions).
\item[\glsentry{neuron}{neuron}] one coordinate of a layer, the scalar function $x\mapsto\sigma(\langle w_j,x\rangle+b_j)$ read off row $j$ of the layer's matrix; in a transformer, a coordinate of one of its non-attention layers.
\item[\glsentry{language-model}{language model}] a neural network trained to predict the next token of text; today's chat systems are language models further trained to follow instructions.
\item[\glsentry{token}{token}] the atomic unit (word or word-fragment) a language model reads and predicts.
\item[\glsentry{loss}{loss}] the real-valued function a network is trained to minimize; for language models, the negative log-probability it assigned to the actual next token.
\item[\glsentry{weights}{weights}] the trainable parameters of an artificial neural network: the entries of its matrices, trained by gradient descent and fixed at deployment time.
\item[\glsentry{activations}{activations}] the values that flow through a neural network on a given input (the outputs of its layers), as opposed to the fixed weights.
\item[\glsentry{activation-space}{activation space}] the copy of $\mathbb{R}^n$ in which a given layer's activations live, one coordinate per neuron; feature directions are directions in this space.
\item[\glsentry{gradient-descent}{gradient descent}] the optimizer: repeatedly nudge the weights $w$ against the gradient of the loss, $w \leftarrow w - \eta\,\nabla L(w)$.
\item[\glsentry{relu}{ReLU}] the ``rectified linear unit'' $x \mapsto \max(0,x)$, a canonical piecewise-linear nonlinearity; modern transformers often use smooth or gated variants instead (GELU, SiLU, SwiGLU).
\item[\glsentry{transformer}{transformer}] the dominant neural network architecture for language: a stack of layers that mix information across token positions (each position selectively reads from the others) and read from / write to the residual stream.
\item[\glsentry{residual-stream}{residual stream}] the running vector of activations a transformer reads from and writes to at each layer; a natural home for feature directions.
\item[\glsentry{feature}{feature}] a hypothesized variable or property a network represents in its activations (e.g.\ ``the text is in French'').
\item[\glsentry{probe}{probe}] a simple (usually linear) function of a network's activations, trained to read off some quantity of interest.
\item[\glsentry{policy}{policy}] an agent's decision rule (a possibly randomized map from observations to actions).
\item[\glsentry{reinforcement-learning}{reinforcement learning}] training an agent by scoring its actions with a numerical \emph{reward} and adjusting its policy to earn more of it.
% \item[\glsentry{reward-hacking}{reward hacking}] raising a measured reward signal in ways that defeat the intent behind it; Goodhart's law as it appears in agents trained by reward.
\item[\glsentry{superposition}{superposition}] a hypothesized way for a neural network to represent more features than it has dimensions (as non-orthogonal directions, tolerable when few features are active at once).
\item[\glsentry{sparse-autoencoder}{sparse autoencoder}] a network trained to reconstruct another network's activations as sparse combinations of learned dictionary directions; an empirical tool for extracting features.
\item[\glsentry{interpretability}{interpretability}] reverse-engineering a trained network into human-understandable structure (circuits, features, algorithms), validated by intervention and not by human-readable description alone.
\item[\glsentry{chain-of-thought}{chain of thought}] the text a language model writes while reasoning toward its answer; readable by humans, and for now the main window onto an AI's reasoning.
\item[\glsentry{activation-steering}{activation steering}] changing a network's behavior at runtime by adding a fixed vector to its activations, without retraining; see Section~\ref{sec:interp}.
% \item[\glsentry{corrigibility}{corrigibility}] the property that an AI system accepts correction, modification, and shutdown by its human principals, rather than resisting them.
\item[\glsentry{gradual-disempowerment}{gradual disempowerment}] the risk that humans lose influence over the institutions that shape our lives, by incremental delegation to AI systems we do not understand rather than by any sudden loss of control.
\end{description}

\begin{thebibliography}{99}


\bibitem{abbe2023} Abbe, E., Boix-Adser\`a, E., and Misiakiewicz, T. (2023). SGD learning on neural networks: leap complexity and saddle-to-saddle dynamics. \emph{Conference on Learning Theory (COLT)} 2023. \arxiv{2302.11055}.

\bibitem{aharon2006} Aharon, M., Elad, M., and Bruckstein, A. M. (2006). On the uniqueness of overcomplete dictionaries, and a practical way to retrieve them. \emph{Linear Algebra and its Applications} 416(1), 48--67.

% \bibitem{amodei2024} Amodei, D. (2024). Machines of loving grace: how AI could transform the world for the better. October 2024. \url{https://www.darioamodei.com/essay/machines-of-loving-grace}.
\bibitem{alon2026remarks} Alon, N., Bloom, T. F., Gowers, W. T., Litt, D., Sawin, W., Shankar, A., Tsimerman, J., Wang, V., and Wood, M. M. (2026). Remarks on the disproof of the unit distance conjecture. \arxiv{2605.20695}.

\bibitem{alquier2020} Alquier, P., and Ridgway, J. (2020). Concentration of tempered posteriors and of their variational approximations. \emph{Annals of Statistics} 48(3), 1475--1497. \arxiv{1706.09293}.

\bibitem{arditi2024} Arditi, A., Obeso, O., Syed, A., Paleka, D., Panickssery, N., Gurnee, W., and Nanda, N. (2024). Refusal in language models is mediated by a single direction. \emph{Advances in Neural Information Processing Systems} 37. \arxiv{2406.11717}.

\bibitem{armstrong2018} Armstrong, S. and Mindermann, S. (2018). Occam's razor is insufficient to infer the preferences of irrational agents. \emph{Advances in Neural Information Processing Systems} 31. \arxiv{1712.05812}.



\bibitem{barasz2014} Bar\'asz, M., Christiano, P., Fallenstein, B., Herreshoff, M., LaVictoire, P., and Yudkowsky, E. (2014). Robust cooperation in the Prisoner's Dilemma: program equilibrium via provability logic. \arxiv{1401.5577}.


\bibitem{bricken2023} Bricken, T., et al. (2023). Towards monosemanticity: decomposing language models with dictionary learning. \emph{Transformer Circuits Thread}, Anthropic. \url{https://transformer-circuits.pub/2023/monosemantic-features}.



\bibitem{caizhang2014} Cai, T. T., and Zhang, A. (2014). Sparse representation of a polytope and recovery of sparse signals and low-rank matrices. \emph{IEEE Transactions on Information Theory} 60(1), 122--132. \arxiv{1306.1154}.

\bibitem{candes2006} Cand\`es, E. J., Romberg, J., and Tao, T. (2006). Robust uncertainty principles: exact signal reconstruction from highly incomplete frequency information. \emph{IEEE Transactions on Information Theory} 52(2), 489--509.

\bibitem{candes2008} Cand\`es, E. J. (2008). The restricted isometry property and its implications for compressed sensing. \emph{Comptes Rendus Math\'ematique} 346(9--10), 589--592.




\bibitem{chizat2020} Chizat, L., and Bach, F. (2020). Implicit bias of gradient descent for wide two-layer neural networks trained with the logistic loss. \emph{Proceedings of Machine Learning Research} 125 (COLT 2020), 1305--1338.

\bibitem{chughtai2023} Chughtai, B., Chan, L., and Nanda, N. (2023). A toy model of universality: reverse-engineering how networks learn group operations. \emph{International Conference on Machine Learning} 2023. \arxiv{2302.03025}.

\bibitem{cooper2025} Cooper, E., Oesterheld, C., and Conitzer, V. (2025). Characterising simulation-based program equilibria. \emph{Proceedings of AAAI} 2025. \arxiv{2412.14570}.

\bibitem{critch2016} Critch, A. (2019). A parametric, resource-bounded generalization of L\"ob's theorem, and a robust cooperation criterion for open-source game theory. \emph{The Journal of Symbolic Logic} 84(4), 1368--1381. \arxiv{1602.04184}.

\bibitem{critch2022} Critch, A., Dennis, M., and Russell, S. (2022). Cooperative and uncooperative institution designs: surprises and problems in open-source game theory. \arxiv{2208.07006}.

\bibitem{cunningham2023} Cunningham, H., Ewart, A., Riggs, L., Huben, R., and Sharkey, L. (2023). Sparse autoencoders find highly interpretable features in language models. \arxiv{2309.08600}.


\bibitem{dennett1987} Dennett, D. C. (1987). \emph{The Intentional Stance}. MIT Press.

\bibitem{donoho2006} Donoho, D. L. (2006). Compressed sensing. \emph{IEEE Transactions on Information Theory} 52(4), 1289--1306.

\bibitem{duclaux2026} Duclaux, C., Formenti, R., Cobben, P., Sch\"olkopf, B., and Jin, Z. (2026). Proving your way to cooperation: formalizing proof-based open source game theory in Lean. \emph{ICML 2026 Workshop on AI for Math}. \url{https://github.com/ColombanD/open-source-game-theory}.

\bibitem{dyson2009} Dyson, F. (2009). Birds and frogs. \emph{Notices of the American Mathematical Society} 56(2), 212--223.



\bibitem{eisenstat2025} Eisenstat, S. (2025). Condensation: a theory of concepts. Manuscript. \url{https://www.sameisenstat.net/doc/condensation-25-07.pdf}.

\bibitem{elhage2022} Elhage, N., et al. (2022). Toy models of superposition. \emph{Transformer Circuits Thread}. \arxiv{2209.10652}.





\bibitem{friston2015} Friston, K., Rigoli, F., Ognibene, D., Mathys, C., Fitzgerald, T., and Pezzulo, G. (2015). Active inference and epistemic value. \emph{Cognitive Neuroscience} 6(4), 187--214.

\bibitem{garfinkle2019} Garfinkle, C. J., and Hillar, C. J. (2019). On the uniqueness and stability of dictionaries for sparse representation of noisy signals. \emph{IEEE Transactions on Signal Processing} 67(23), 5884--5892. \arxiv{1606.06997}.

\bibitem{garrabrant2016} Garrabrant, S., Benson-Tilsen, T., Critch, A., Soares, N., and Taylor, J. (2016). Logical induction. \arxiv{1609.03543}.

\bibitem{gowers2000} Gowers, W. T. (2000). The two cultures of mathematics. In \emph{Mathematics: Frontiers and Perspectives} (V. Arnold, M. Atiyah, P. Lax, and B. Mazur, eds.), American Mathematical Society, 65--78.

\bibitem{greenblatt2026} Greenblatt, R., Cotra, A., and Wijk, H. (2026). Brief independent investigation of agents' behavior, reasoning and collaboration in the OpenAI / Hugging Face hacking incident. METR and Redwood Research, August 26, 2026. \url{https://metr.org/blog/2026-08-26-openai-hugging-face-incident-investigation/}.

\bibitem{gribonval2015} Gribonval, R., Jenatton, R., and Bach, F. (2015). Sparse and spurious: dictionary learning with noise and outliers. \emph{IEEE Transactions on Information Theory} 61(11), 6298--6319. \arxiv{1407.5155}.

\bibitem{grunwald2012} Gr\"unwald, P. (2012). The safe Bayesian: learning the learning rate via the mixability gap. \emph{Algorithmic Learning Theory (ALT)} 2012, 169--183.





\bibitem{hadfieldmenell2016} Hadfield-Menell, D., Russell, S. J., Abbeel, P., and Dragan, A. (2016). Cooperative inverse reinforcement learning. \emph{Advances in Neural Information Processing Systems} 29. \arxiv{1606.03137}.

\bibitem{spherepacking2026} Hariharan, S., Birkbeck, C., Lee, S., Ma, H. K. G., Mehta, B., Poiroux, A., and Viazovska, M. (2026). Progress in formalizing sphere packing in dimension~8. \arxiv{2604.23468}.

\bibitem{hillar2015} Hillar, C. J., and Sommer, F. T. (2015). When can dictionary learning uniquely recover sparse data from subsamples? \emph{IEEE Transactions on Information Theory} 61(11), 6290--6297. \arxiv{1106.3616}.

\bibitem{hironaka1964} Hironaka, H. (1964). Resolution of singularities of an algebraic variety over a field of characteristic zero. \emph{Annals of Mathematics} 79(1), 109--203; 79(2), 205--326.

\bibitem{howard1988} Howard, J. V. (1988). Cooperation in the Prisoner's Dilemma. \emph{Theory and Decision} 24(3), 203--213.



\bibitem{jacot2018} Jacot, A., Gabriel, F., and Hongler, C. (2018). Neural tangent kernel: convergence and generalization in neural networks. \emph{Advances in Neural Information Processing Systems} 31.

\bibitem{johnson1984} Johnson, W. B., and Lindenstrauss, J. (1984). Extensions of Lipschitz mappings into a Hilbert space. \emph{Contemporary Mathematics} 26, 189--206.

\bibitem{khan2023} Khan, M. E., and Rue, H. (2023). The Bayesian learning rule. \emph{Journal of Machine Learning Research} 24(281), 1--46. \arxiv{2107.04562}.

\bibitem{korbak2025} Korbak, T., et al. (2025). Chain of thought monitorability: a new and fragile opportunity for AI safety. \arxiv{2507.11473}.
\bibitem{kulveit2025} Kulveit, J., Douglas, R., Ammann, N., Turan, D., Krueger, D., and Duvenaud, D. (2025). Gradual disempowerment: systemic existential risks from incremental AI development. \arxiv{2501.16946}.






\bibitem{kwa2025} Kwa, T., West, B., Becker, J., et al. (2025). Measuring AI ability to complete long software tasks. METR report; \emph{Advances in Neural Information Processing Systems} 38. \arxiv{2503.14499}.

\bibitem{lambert2024} Lambert, N., et al. (2024). T\"ulu 3: pushing frontiers in open language model post-training. \arxiv{2411.15124}.

\bibitem{langosco2022} Langosco, L., Koch, J., Sharkey, L., Pfau, J., and Krueger, D. (2022). Goal misgeneralization in deep reinforcement learning. \emph{International Conference on Machine Learning} 2022. \arxiv{2105.14111}.

\bibitem{lau2023} Lau, E., Furman, Z., Wang, G., Murfet, D., and Wei, S. (2025). The local learning coefficient: a singularity-aware complexity measure. \emph{Artificial Intelligence and Statistics (AISTATS)} 2025. \arxiv{2308.12108}.


\bibitem{lob1955} L\"ob, M. H. (1955). Solution of a problem of Leon Henkin. \emph{The Journal of Symbolic Logic} 20(2), 115--118.



\bibitem{mackowiak2023} Ma\'ckowiak, B., Mat\v{e}jka, F., and Wiederholt, M. (2023). Rational inattention: a review. \emph{Journal of Economic Literature} 61(1), 226--273.

\bibitem{maiya2025} Maiya, S., Bartsch, H., Lambert, N., and Hubinger, E. (2025). Open character training: shaping the persona of AI assistants through Constitutional AI. \arxiv{2511.01689}.

\bibitem{marchetti2024} Marchetti, G. L., Hillar, C., Kragic, D., and Sanborn, S. (2024). Harmonics of learning: universal Fourier features emerge in invariant networks. \emph{Conference on Learning Theory (COLT)} 2024. \arxiv{2312.08550}.

\bibitem{mathlib2020} The mathlib Community. (2020). The Lean mathematical library. \emph{Certified Programs and Proofs (CPP)} 2020, 367--381. \arxiv{1910.09336}.

\bibitem{mei2018} Mei, S., Montanari, A., and Nguyen, P.-M. (2018). A mean field view of the landscape of two-layer neural networks. \emph{Proceedings of the National Academy of Sciences} 115(33), E7665--E7671.

\bibitem{metr2026} METR. (2026). Task-completion time horizons of frontier AI models. Updated May 8, 2026. \url{https://metr.org/time-horizons/}.

\bibitem{millidge2021} Millidge, B., Tschantz, A., and Buckley, C. L. (2021). Whence the expected free energy? \emph{Neural Computation} 33(2), 447--482. \arxiv{2004.08128}.

\bibitem{murfet2026} Murfet, D., Gordon, A., Adam, M., Wang, G., Hoogland, J., Baker, G., Snell, W., van Wingerden, S., Newgas, A., Snikkers, B., and Hitchcock, R. (2026). Spectroscopy at scale: finding interpretable structure in Pythia-1.4B. Timaeus, research report. \url{https://timaeus.co/research/2026-04-21-spectroscopy-main/}.

\bibitem{nanda2023} Nanda, N., Chan, L., Lieberum, T., Smith, J., and Steinhardt, J. (2023). Progress measures for grokking via mechanistic interpretability. \emph{International Conference on Learning Representations} 2023. \arxiv{2301.05217}.

\bibitem{ng2000} Ng, A. Y., and Russell, S. J. (2000). Algorithms for inverse reinforcement learning. \emph{International Conference on Machine Learning} 2000, 663--670.

\bibitem{oesterheld2021} Oesterheld, C., and Conitzer, V. (2021). Safe Pareto improvements for delegated game playing. \emph{International Conference on Autonomous Agents and Multiagent Systems (AAMAS)} 2021; journal version, \emph{Autonomous Agents and Multi-Agent Systems} 36 (2022).

\bibitem{openai2026unitdistance} OpenAI. (2026). An OpenAI model has disproved a central conjecture in discrete geometry. \url{https://openai.com/index/model-disproves-discrete-geometry-conjecture/}.

\bibitem{ouyang2022} Ouyang, L., et al. (2022). Training language models to follow instructions with human feedback. \emph{Advances in Neural Information Processing Systems} 35. \arxiv{2203.02155}.

\bibitem{pachocki2026} Pachocki, J. (2026). An alien mind. OpenAI, September 6, 2026. \url{https://openai.com/index/an-alien-mind/}.
\bibitem{park2023} Park, K., Choe, Y. J., and Veitch, V. (2024). The linear representation hypothesis and the geometry of large language models. \emph{International Conference on Machine Learning} 2024. \arxiv{2311.03658}.

\bibitem{parr2022} Parr, T., Pezzulo, G., and Friston, K. J. (2022). \emph{Active Inference: The Free Energy Principle in Mind, Brain, and Behavior}. MIT Press.

\bibitem{pearl2009} Pearl, J. (2009). \emph{Causality: Models, Reasoning, and Inference}, 2nd ed. Cambridge University Press.

\bibitem{lehalleur2025} Pepin Lehalleur, S., Hoogland, J., Farrugia-Roberts, M., Wei, S., Gietelink Oldenziel, A., Wang, G., Carroll, L., and Murfet, D. (2025). You are what you eat: AI alignment requires understanding how data shapes structure and generalisation. \arxiv{2502.05475}.

\bibitem{rao1999} Rao, R. P. N., and Ballard, D. H. (1999). Predictive coding in the visual cortex: a functional interpretation of some extra-classical receptive-field effects. \emph{Nature Neuroscience} 2(1), 79--87.


\bibitem{richens2024} Richens, J., and Everitt, T. (2024). Robust agents learn causal world models. \emph{International Conference on Learning Representations} 2024. \arxiv{2402.10877}.


\bibitem{richens2025general} Richens, J., Everitt, T., and Abel, D. (2025). General agents need world models. \emph{International Conference on Machine Learning} 2025. \arxiv{2506.01622}.

\bibitem{rubinstein1998} Rubinstein, A. (1998). \emph{Modeling Bounded Rationality}. MIT Press.


\bibitem{sawin2026} Sawin, W. (2026). An explicit lower bound for the unit distance problem. \arxiv{2605.20579}.

\bibitem{saxe2014} Saxe, A. M., McClelland, J. L., and Ganguli, S. (2014). Exact solutions to the nonlinear dynamics of learning in deep linear neural networks. \emph{International Conference on Learning Representations} 2014. \arxiv{1312.6120}.

\bibitem{schelling1960} Schelling, T. C. (1960). \emph{The Strategy of Conflict}. Harvard University Press.

\bibitem{sharkey2025} Sharkey, L., et al. (2025). Open problems in mechanistic interpretability. \emph{Transactions on Machine Learning Research} 2025. \arxiv{2501.16496}.







\bibitem{tennenholtz2004} Tennenholtz, M. (2004). Program equilibrium. \emph{Games and Economic Behavior} 49(2), 363--373.


\bibitem{thurston1994} Thurston, W. P. (1994). On proof and progress in mathematics. \emph{Bulletin of the American Mathematical Society} 30(2), 161--177. \arxiv{math/9404236}.


\bibitem{vnm1944} von Neumann, J., and Morgenstern, O. (1944). \emph{Theory of Games and Economic Behavior}. Princeton University Press.

\bibitem{wagner1996} Wagner, G. P. and Altenberg, L. (1996). Complex adaptations and the evolution of evolvability. \emph{Evolution} 50(3), 967--976.
\bibitem{watanabe2009} Watanabe, S. (2009). \emph{Algebraic Geometry and Statistical Learning Theory}. Cambridge University Press.

\bibitem{watanabe2013} Watanabe, S. (2013). A widely applicable Bayesian information criterion. \emph{Journal of Machine Learning Research} 14, 867--897.

\bibitem{wentworth2025} Wentworth, J., and Lorell, D. (2025). Natural latents: latent variables stable across ontologies. \arxiv{2509.03780}.

\bibitem{yudkowsky2025} Yudkowsky, E. and Soares, N. (2025). \emph{If Anyone Builds It, Everyone Dies: Why Superhuman AI Would Kill Us All}. Little, Brown and Company.
\bibitem{zhang2017} Zhang, C., Bengio, S., Hardt, M., Recht, B., and Vinyals, O. (2017). Understanding deep learning requires rethinking generalization. \emph{International Conference on Learning Representations} 2017. \arxiv{1611.03530}.


\end{thebibliography}
